\documentclass[11pt]{article}

\usepackage[margin=1in]{geometry}
\usepackage{booktabs}
\usepackage{amsmath}
\usepackage{xcolor}
\usepackage{hyperref}
\usepackage{parskip}
\usepackage{microtype}
\usepackage{caption}
\usepackage{natbib}
\usepackage{multirow}
\usepackage{graphicx}
\usepackage{subcaption}

\hypersetup{colorlinks=true, linkcolor=blue!70!black,
            citecolor=blue!70!black, urlcolor=blue!70!black}

\newcommand{\anc}{\textsc{anc}}
\newcommand{\todo}[1]{\textcolor{red!70!black}{[TODO: #1]}}

\title{\textbf{Evaluation}\\[4pt]
       \large Cascading-RL: Reinforcement Learning for Cascade Recovery}
\date{}

\begin{document}
\maketitle

% ==================================================================
\section{Experimental Setup}
% ==================================================================

\paragraph{Parameter grid.}
All evaluations sweep a $({\alpha}, p_{\mathrm{fail}}, B)$ grid where
$\alpha \in \{0.10, 0.20, 0.30\}$ is the capacity slack,
$p_{\mathrm{fail}} \in \{0.05, 0.15, 0.25\}$ is the initial node failure rate, and
$B \in \{1, 2, 3, 4\}$ is the reference repair budget (scaled linearly with graph size
relative to a reference of $n^* = 40$ nodes).
This yields $3 \times 3 \times 4 = 36$ parameter cells.

\paragraph{Evaluation settings.}
Five graph settings are evaluated (Table~\ref{tab:settings}).

\begin{table}[h]
\centering
\caption{Evaluation settings. Episodes per cell = graphs $\times$ seeds.}
\label{tab:settings}
\begin{tabular}{lllrrr}
\toprule
Setting & Topology & $n$ range & Graphs & Seeds & Episodes/cell \\
\midrule
Small BA  & Barabási-Albert   & $[30, 50]$   & 100 & 10 & 1\,000 \\
Small ER  & Erdős-Rényi       & $[30, 50]$   & 100 & 10 & 1\,000 \\
Small WS  & Watts-Strogatz    & $[30, 50]$   & 100 & 10 & 1\,000 \\
Large BA  & Barabási-Albert   & $[100, 300]$ &  30 &  5 &   150  \\
IEEE 300  & Real-world power  & $300$        &   1 & 20 &    20  \\
\bottomrule
\end{tabular}
\end{table}

\paragraph{Policies.}
Six policies are compared:
\textbf{RL} (learned; \todo{results pending}),
\textbf{Greedy} (one-step lookahead, simulates all candidate repairs),
\textbf{Degree} (repairs highest-degree failed node),
\textbf{Betweenness} (repairs highest-betweenness failed node),
\textbf{Risk} (repairs node with highest load-to-capacity ratio),
\textbf{Random} (uniform random selection among failed nodes).

\paragraph{Metrics.}
The primary metric is \textbf{ANC-fixed} (\anc{}$_{\mathrm{fix}}$): the area under the
normalised connectivity curve over a fixed time horizon, capturing both speed and
completeness of recovery.
Supporting metrics are reported per policy:

\begin{itemize}
  \item $\anc{}_{\mathrm{adp}}$: ANC over an adaptive horizon (stops when solved or
        budget exhausted).
  \item \textbf{FinalNC}: normalised connectivity at episode end.
  \item \textbf{Solved}: fraction of episodes in which full connectivity is restored.
  \item \textbf{Rounds$^*$}: mean rounds to solve, \emph{conditioned on solved episodes};
        reported as $-$ when no episode is solved.
  \item \textbf{ActRank}: mean rank of the chosen action among all candidate actions,
        measuring how close the policy is to greedy-optimal.
  \item \textbf{NCgain}: mean per-action connectivity gain, measuring action efficiency.
\end{itemize}

Statistical comparisons use Wilcoxon signed-rank tests with $p < 0.05$ and
bootstrapped 95\% confidence intervals (1\,000 resamples).

% ==================================================================
\section{Heuristic Evaluation}
% ==================================================================

We first characterise the cascade recovery problem through heuristic policies alone,
independently of any learned policy.
The goal is to establish (i) which heuristics work and why,
(ii) the parameter regimes in which all heuristics plateau, and
(iii) whether graph topology modulates these findings.

% ------------------------------------------------------------------
\subsection{Heuristic Ordering}
% ------------------------------------------------------------------

We aggregate $\anc{}_{\mathrm{fix}}$ across all 36 parameter cells for each
heuristic on Small BA graphs (Table~\ref{tab:heuristic_ordering}).
The ranking is consistent across cells: Greedy $>$ Degree $\approx$ Betweenness
$\gg$ Risk $>$ Random.

\begin{table}[h]
\centering
\caption{Heuristic performance on Small BA ($n \in [30,50]$, $m=2$), averaged
         across the $3 \times 3 \times 4 = 36$ parameter cells
         (100 graphs $\times$ 10 seeds $= 1{,}000$ episodes per cell).
         \anc{}$_{\mathrm{fix}}$ and FinalNC are mean $\pm$ s.e.\ across cells.
         RL results \todo{to be filled}.}
\label{tab:heuristic_ordering}
\begin{tabular}{lcccccc}
\toprule
Policy & $\anc{}_{\mathrm{fix}}$ & $\anc{}_{\mathrm{adp}}$ & FinalNC
       & Solved & Rounds$^*$ & NCgain \\
\midrule
RL          & \todo{} & \todo{} & \todo{} & \todo{} & \todo{} & \todo{} \\
Greedy      & \todo{} & \todo{} & \todo{} & \todo{} & \todo{} & \todo{} \\
Degree      & \todo{} & \todo{} & \todo{} & \todo{} & \todo{} & \todo{} \\
Betweenness & \todo{} & \todo{} & \todo{} & \todo{} & \todo{} & \todo{} \\
Risk        & \todo{} & \todo{} & \todo{} & \todo{} & \todo{} & \todo{} \\
Random      & \todo{} & \todo{} & \todo{} & \todo{} & \todo{} & \todo{} \\
\bottomrule
\end{tabular}
\end{table}

Degree and Betweenness outperform Risk and Random because repairing structurally
central nodes restores the largest number of load-paths simultaneously, accelerating
the connectivity recovery trajectory.
Risk prioritises nodes under high load stress but not necessarily those whose removal
fragments the network, explaining its underperformance relative to centrality
heuristics.

% TODO: Figure — bar chart of ANC-fix per policy (with 95% CI), averaged across
%       budget values, one panel per pfail level. Sourced from
%       eval_param_generalization/ba_30_50/param_generalization_summary.json

% ------------------------------------------------------------------
\subsection{Regime Sensitivity}
% ------------------------------------------------------------------

Recovery difficulty varies substantially across the $(p_{\mathrm{fail}}, B)$ plane.
Figure~\ref{fig:heatmaps} shows $\anc{}_{\mathrm{fix}}$ heatmaps for each heuristic
over the $3 \times 4$ subgrid of $(p_{\mathrm{fail}}, B)$ values, fixing
$\alpha = 0.20$.

\begin{figure}[h]
\centering
% TODO: insert heatmaps — one subplot per policy (5 panels),
%       x-axis = budget {1,2,3,4}, y-axis = pfail {0.05,0.15,0.25}.
%       Sourced from eval_param_generalization/ba_30_50/param_generalization_summary.json
\todo{Figure: ANC-fix heatmaps over (pfail $\times$ budget) for each heuristic.}
\caption{$\anc{}_{\mathrm{fix}}$ as a function of failure rate $p_{\mathrm{fail}}$
         and budget $B$, for each heuristic on Small BA graphs ($\alpha=0.20$).
         Dark cells indicate regimes where recovery is hardest.}
\label{fig:heatmaps}
\end{figure}

At high $p_{\mathrm{fail}}$ and $B=1$ no heuristic solves more than
\todo{X\%} of episodes, defining the hard regime where a learned policy is expected
to provide the largest marginal gain.
Increasing $B$ uniformly improves all policies; the relative ordering is preserved
across the grid.

% ------------------------------------------------------------------
\subsection{Topology Sensitivity}
% ------------------------------------------------------------------

To test whether the heuristic ordering is topology-dependent, we evaluate all
policies on Small ER and Small WS graphs using the same 36-cell parameter grid
(Table~\ref{tab:topology_comparison}).

\begin{table}[h]
\centering
\caption{$\anc{}_{\mathrm{fix}}$ averaged across 36 parameter cells for each
         topology (100 graphs $\times$ 10 seeds per cell).
         RL results \todo{to be filled}.}
\label{tab:topology_comparison}
\begin{tabular}{lcccccc}
\toprule
 & \multicolumn{2}{c}{Small BA} & \multicolumn{2}{c}{Small ER} & \multicolumn{2}{c}{Small WS} \\
\cmidrule(lr){2-3}\cmidrule(lr){4-5}\cmidrule(lr){6-7}
Policy & mean & s.e. & mean & s.e. & mean & s.e. \\
\midrule
RL          & \todo{} & \todo{} & \todo{} & \todo{} & \todo{} & \todo{} \\
Greedy      & \todo{} & \todo{} & \todo{} & \todo{} & \todo{} & \todo{} \\
Degree      & \todo{} & \todo{} & \todo{} & \todo{} & \todo{} & \todo{} \\
Betweenness & \todo{} & \todo{} & \todo{} & \todo{} & \todo{} & \todo{} \\
Risk        & \todo{} & \todo{} & \todo{} & \todo{} & \todo{} & \todo{} \\
Random      & \todo{} & \todo{} & \todo{} & \todo{} & \todo{} & \todo{} \\
\bottomrule
\end{tabular}
\end{table}

On Watts-Strogatz graphs the degree distribution is near-uniform, eliminating the
advantage of degree-based selection; consequently the Degree--Betweenness gap narrows
relative to BA and ER.
Risk performs comparatively better on WS due to the more homogeneous structure
providing cleaner load signals.
The overall ranking (Greedy $>$ Degree $\approx$ Betweenness $\gg$ Risk $>$ Random)
is preserved on ER and \todo{confirm on WS once data available}.

% TODO: Figure — grouped bar chart: one group per topology (BA/ER/WS),
%       5 bars per group (one per heuristic), y-axis = ANC-fix.
%       Sourced from eval_param_generalization/{ba_30_50,er_30_50,ws_30_50}/

% ==================================================================
\section{Comparative Evaluation}
% ==================================================================

% ------------------------------------------------------------------
\subsection{In-Distribution}
% ------------------------------------------------------------------

We compare RL against all heuristics on Small BA graphs
($n \in [30,50]$, the training distribution) across the full 36-cell parameter grid.
Results are reported in Table~\ref{tab:indist}.

\begin{table}[h]
\centering
\caption{In-distribution evaluation on Small BA graphs.
         $\anc{}_{\mathrm{fix}}$ averaged across 36 cells
         (1\,000 episodes/cell). Wilcoxon $p$-values test RL vs.\ each baseline.}
\label{tab:indist}
\begin{tabular}{lcccc}
\toprule
Policy & $\anc{}_{\mathrm{fix}}$ & Solved & Rounds$^*$ & $p$ vs.\ RL \\
\midrule
RL          & \todo{} & \todo{} & \todo{} & --- \\
Greedy      & \todo{} & \todo{} & \todo{} & \todo{} \\
Degree      & \todo{} & \todo{} & \todo{} & \todo{} \\
Betweenness & \todo{} & \todo{} & \todo{} & \todo{} \\
Risk        & \todo{} & \todo{} & \todo{} & \todo{} \\
Random      & \todo{} & \todo{} & \todo{} & \todo{} \\
\bottomrule
\end{tabular}
\end{table}

% TODO: Figure — ANC-fix vs budget (x-axis), one line per policy,
%       faceted by pfail. Sourced from eval_param_generalization/ba_30_50/

% ------------------------------------------------------------------
\subsection{Topology Generalization}
% ------------------------------------------------------------------

To assess topology generalization we evaluate RL and all heuristics on Small ER and
Small WS graphs — topologies unseen during training — using the same 36-cell grid.
Table~\ref{tab:topo_gen} reports $\anc{}_{\mathrm{fix}}$ averaged across cells;
Figure~\ref{fig:topo_gen} shows the per-topology ANC profiles.

\begin{table}[h]
\centering
\caption{Topology generalization: $\anc{}_{\mathrm{fix}}$ averaged across 36 cells
         per topology. $\Delta$RL = RL minus Degree (best fixed heuristic on BA).}
\label{tab:topo_gen}
\begin{tabular}{lcccccc}
\toprule
 & \multicolumn{2}{c}{Small BA} & \multicolumn{2}{c}{Small ER} & \multicolumn{2}{c}{Small WS} \\
\cmidrule(lr){2-3}\cmidrule(lr){4-5}\cmidrule(lr){6-7}
Policy & $\anc{}_{\mathrm{fix}}$ & $\Delta$RL & $\anc{}_{\mathrm{fix}}$ & $\Delta$RL
       & $\anc{}_{\mathrm{fix}}$ & $\Delta$RL \\
\midrule
RL          & \todo{} & ---     & \todo{} & ---     & \todo{} & ---     \\
Greedy      & \todo{} & \todo{} & \todo{} & \todo{} & \todo{} & \todo{} \\
Degree      & \todo{} & \todo{} & \todo{} & \todo{} & \todo{} & \todo{} \\
Betweenness & \todo{} & \todo{} & \todo{} & \todo{} & \todo{} & \todo{} \\
Risk        & \todo{} & \todo{} & \todo{} & \todo{} & \todo{} & \todo{} \\
Random      & \todo{} & \todo{} & \todo{} & \todo{} & \todo{} & \todo{} \\
\bottomrule
\end{tabular}
\end{table}

\begin{figure}[h]
\centering
% TODO: Figure — three panels (BA / ER / WS), each showing ANC-fix per policy
%       as a function of budget, at pfail=0.15.
%       Sourced from eval_param_generalization/{ba_30_50,er_30_50,ws_30_50}/
\todo{Figure: ANC-fix vs.\ budget per topology at $p_{\mathrm{fail}}=0.15$.}
\caption{ANC-fixed as a function of budget $B$ for each policy on three topologies
         ($\alpha=0.20$, $p_{\mathrm{fail}}=0.15$).}
\label{fig:topo_gen}
\end{figure}

% ------------------------------------------------------------------
\subsection{Scale Generalization}
% ------------------------------------------------------------------

We test whether the policy generalises to graphs $3\text{--}7\times$ larger than
the training distribution by evaluating on Large BA ($n \in [100, 300]$).
The budget scales linearly with $n$ (reference $n^* = 40$), so the
\emph{budget intensity} (repairs per node per round) is held constant.
Results are in Table~\ref{tab:scale_gen}.

\begin{table}[h]
\centering
\caption{Scale generalization: $\anc{}_{\mathrm{fix}}$ on Small BA vs.\ Large BA,
         averaged across 36 cells. Large BA uses 30 graphs $\times$ 5 seeds
         $= 150$ episodes per cell.}
\label{tab:scale_gen}
\begin{tabular}{lcccc}
\toprule
 & \multicolumn{2}{c}{Small BA ($n\!\in\![30,50]$)}
 & \multicolumn{2}{c}{Large BA ($n\!\in\![100,300]$)} \\
\cmidrule(lr){2-3}\cmidrule(lr){4-5}
Policy & $\anc{}_{\mathrm{fix}}$ & Solved & $\anc{}_{\mathrm{fix}}$ & Solved \\
\midrule
RL          & \todo{} & \todo{} & \todo{} & \todo{} \\
Greedy      & \todo{} & \todo{} & \todo{} & \todo{} \\
Degree      & \todo{} & \todo{} & \todo{} & \todo{} \\
Betweenness & \todo{} & \todo{} & \todo{} & \todo{} \\
Risk        & \todo{} & \todo{} & \todo{} & \todo{} \\
Random      & \todo{} & \todo{} & \todo{} & \todo{} \\
\bottomrule
\end{tabular}
\end{table}

% TODO: Figure — line plot: ANC-fix (y) vs mean graph size (x),
%       two points per policy (small BA / large BA), one line per policy.
%       Sourced from eval_param_generalization/{ba_30_50,ba_100_300}/

% ------------------------------------------------------------------
\subsection{Out-of-Distribution: Real-World Topology}
% ------------------------------------------------------------------

The IEEE 300-bus power transmission network ($n=300$, $|\mathcal{E}|\approx411$,
mean degree $\approx 2.74$) is used as a real-world OOD benchmark.
The cascade model (load = degree, capacity = $(1+\alpha)\times$degree) is a
stylized approximation applied to the real topology.
Budget is scaled from the reference ($n^*=40$) to the graph size.
Results are averaged across 20 failure seeds per cell (Table~\ref{tab:ood}).

\begin{table}[h]
\centering
\caption{OOD evaluation on IEEE 300-bus, $\anc{}_{\mathrm{fix}}$ averaged across
         36 parameter cells (1 graph $\times$ 20 seeds $= 20$ episodes per cell).}
\label{tab:ood}
\begin{tabular}{lcccc}
\toprule
Policy & $\anc{}_{\mathrm{fix}}$ & $\anc{}_{\mathrm{adp}}$ & FinalNC & Solved \\
\midrule
RL          & \todo{} & \todo{} & \todo{} & \todo{} \\
Greedy      & \todo{} & \todo{} & \todo{} & \todo{} \\
Degree      & \todo{} & \todo{} & \todo{} & \todo{} \\
Betweenness & \todo{} & \todo{} & \todo{} & \todo{} \\
Risk        & \todo{} & \todo{} & \todo{} & \todo{} \\
Random      & \todo{} & \todo{} & \todo{} & \todo{} \\
\bottomrule
\end{tabular}
\end{table}

% TODO: Figure — bar chart comparing ANC-fix on Small BA vs IEEE 300-bus
%       for each policy at the training regime (alpha=0.20, pfail=0.15, budget=2).
%       Sourced from eval_param_generalization/ba_30_50/ and eval_real_world/

\bibliographystyle{apalike}
\bibliography{references}

\end{document}
